\section{AIME I 1998}
        \begin{enumerate}
        	\item (\textit{AIME I 1998}) For how many values of $k$ is $12^{12}$ the \href{https://artofproblemsolving.com/wiki/index.php?title=Least\_common\_multiple}{least common multiple} of the positive integers $6^6$, $8^8$, and $k$?


 

	\item (\textit{AIME I 1998}) Find the number of \href{https://artofproblemsolving.com/wiki/index.php?title=Ordered\_pair}{ordered pairs} $(x,y)$ of positive integers that satisfy $x \le 2y \le 60$ and $y \le 2x \le 60$.


 

	\item (\textit{AIME I 1998}) The graph of $y^2 + 2xy + 40|x|= 400$ partitions the plane into several regions.  What is the area of the bounded region?


 

	\item (\textit{AIME I 1998}) Nine tiles are numbered $1, 2, 3, \cdots, 9,$ respectively.  Each of three players randomly selects and keeps three of the tiles, and sums those three values.  The \href{https://artofproblemsolving.com/wiki/index.php?title=Probability}{probability} that all three players obtain an \href{https://artofproblemsolving.com/wiki/index.php?title=Odd}{odd} sum is $m/n,$ where $m$ and $n$ are \href{https://artofproblemsolving.com/wiki/index.php?title=Relatively\_prime}{relatively prime} \href{https://artofproblemsolving.com/wiki/index.php?title=Positive\_integer}{positive integers}.  Find $m+n.$


	\item (\textit{AIME I 1998}) Given that $A_k = \frac {k(k - 1)}2\cos\frac {k(k - 1)\pi}2,$ find $\left\lvert A_{19} + A_{20} + \cdots + A_{98}\right\rvert.$


	\item (\textit{AIME I 1998}) Let $ABCD$ be a \href{https://artofproblemsolving.com/wiki/index.php?title=Parallelogram}{parallelogram}.  Extend $\overline{DA}$ through $A$ to a point $P,$ and let $\overline{PC}$ meet $\overline{AB}$ at $Q$ and $\overline{DB}$ at $R.$  Given that $PQ = 735$ and $QR = 112,$ find $RC.$


	\item (\textit{AIME I 1998}) Let $n$ be the number of ordered quadruples $(x_1,x_2,x_3,x_4)$ of positive odd \href{https://artofproblemsolving.com/wiki/index.php?title=Integer}{integers} that satisfy $\sum_{i = 1}^4 x_i = 98.$  Find $\frac n{100}.$


	\item (\textit{AIME I 1998}) Except for the first two terms, each term of the sequence $1000, x, 1000 - x,\ldots$ is obtained by subtracting the preceding term from the one before that.  The last term of the sequence is the first \href{https://artofproblemsolving.com/wiki/index.php?title=Negative}{negative} term encountered.  What positive integer $x$ produces a sequence of maximum length?


 

	\item (\textit{AIME I 1998}) Two mathematicians take a morning coffee break each day.  They arrive at the cafeteria independently, at random times between 9 a.m. and 10 a.m., and stay for exactly $m$ minutes.  The \href{https://artofproblemsolving.com/wiki/index.php?title=Probability}{probability} that either one arrives while the other is in the cafeteria is $40 \%,$ and $m = a - b\sqrt {c},$ where $a, b,$ and $c$ are \href{https://artofproblemsolving.com/wiki/index.php?title=Positive}{positive} \href{https://artofproblemsolving.com/wiki/index.php?title=Integer}{integers}, and $c$ is not divisible by the square of any \href{https://artofproblemsolving.com/wiki/index.php?title=Prime}{prime}.  Find $a + b + c.$


	\item (\textit{AIME I 1998}) Eight \href{https://artofproblemsolving.com/wiki/index.php?title=Sphere}{spheres} of \href{https://artofproblemsolving.com/wiki/index.php?title=Radius}{radius} 100 are placed on a flat \href{https://artofproblemsolving.com/wiki/index.php?title=Plane}{surface} so that each sphere is \href{https://artofproblemsolving.com/wiki/index.php?title=Tangent}{tangent} to two others and their \href{https://artofproblemsolving.com/wiki/index.php?title=Center}{centers} are the vertices of a regular \href{https://artofproblemsolving.com/wiki/index.php?title=Octagon}{octagon}.  A ninth sphere is placed on the flat surface so that it is tangent to each of the other eight spheres.  The radius of this last sphere is $a +b\sqrt {c},$ where $a, b,$ and $c$ are \href{https://artofproblemsolving.com/wiki/index.php?title=Positive}{positive} \href{https://artofproblemsolving.com/wiki/index.php?title=Integer}{integers}, and $c$ is not divisible by the square of any \href{https://artofproblemsolving.com/wiki/index.php?title=Prime}{prime}.  Find $a + b + c$.


 

	\item (\textit{AIME I 1998}) Three of the edges of a \href{https://artofproblemsolving.com/wiki/index.php?title=Cube}{cube} are $\overline{AB}, \overline{BC},$ and $\overline{CD},$ and $\overline{AD}$ is an interior \href{https://artofproblemsolving.com/wiki/index.php?title=Diagonal}{diagonal}.  Points $P, Q,$ and $R$ are on $\overline{AB}, \overline{BC},$ and $\overline{CD},$ respectively, so that $AP = 5, PB = 15, BQ = 15,$ and $CR = 10.$  What is the \href{https://artofproblemsolving.com/wiki/index.php?title=Area}{area} of the \href{https://artofproblemsolving.com/wiki/index.php?title=Polygon}{polygon} that is the \href{https://artofproblemsolving.com/wiki/index.php?title=Intersection}{intersection} of \href{https://artofproblemsolving.com/wiki/index.php?title=Plane}{plane} $PQR$ and the cube?


 

	\item (\textit{AIME I 1998}) Let $ABC$ be \href{https://artofproblemsolving.com/wiki/index.php?title=Equilateral\_triangle}{equilateral}, and $D, E,$ and $F$ be the \href{https://artofproblemsolving.com/wiki/index.php?title=Midpoint}{midpoints} of $\overline{BC}, \overline{CA},$ and $\overline{AB},$ respectively.  There exist \href{https://artofproblemsolving.com/wiki/index.php?title=Point}{points} $P, Q,$ and $R$ on $\overline{DE}, \overline{EF},$ and $\overline{FD},$ respectively, with the property that $P$ is on $\overline{CQ}, Q$ is on $\overline{AR},$ and $R$ is on $\overline{BP}.$  The \href{https://artofproblemsolving.com/wiki/index.php?title=Ratio}{ratio} of the area of triangle $ABC$ to the area of triangle $PQR$ is $a + b\sqrt {c},$ where $a, b,$ and $c$ are integers, and $c$ is not divisible by the square of any \href{https://artofproblemsolving.com/wiki/index.php?title=Prime}{prime}.  What is $a^{2} + b^{2} + c^{2}$?


 

	\item (\textit{AIME I 1998}) If $\{a_1,a_2,a_3,\ldots,a_n\}$ is a \href{https://artofproblemsolving.com/wiki/index.php?title=Set}{set} of \href{https://artofproblemsolving.com/wiki/index.php?title=Real\_numbers}{real numbers}, indexed so that $a_1 < a_2 < a_3 < \cdots < a_n,$ its \textit{complex power sum} is defined to be $a_1i + a_2i^2+ a_3i^3 + \cdots + a_ni^n,$ where $i^2 = - 1.$  Let $S_n$ be the sum of the complex power sums of all nonempty \href{https://artofproblemsolving.com/wiki/index.php?title=Subset}{subsets} of $\{1,2,\ldots,n\}.$  Given that $S_8 = - 176 - 64i$ and $S_9 = p + qi,$ where $p$ and $q$ are integers, find $\left\lvert p\right\rvert + \left\lvert q\right\rvert.$


	\item (\textit{AIME I 1998}) An $m\times n\times p$ rectangular box has half the volume of an $(m + 2)\times(n + 2)\times(p + 2)$ rectangular box, where $m, n,$ and $p$ are integers, and $m\le n\le p.$  What is the largest possible value of $p$?


 

	\item (\textit{AIME I 1998}) Define a \textit{domino} to be an ordered pair of \textit{distinct} positive integers. A \textit{proper sequence} of dominos is a list of distinct dominos in which the first coordinate of each pair after the first equals the second coordinate of the immediately preceding pair, and in which $(i,j)$ and $(j,i)$ do not \textit{both} appear for any $i$ and $j$.  Let $D_{40}$ be the set of all dominos whose coordinates are no larger than 40.  Find the length of the longest proper sequence of dominos that can be formed using the dominos of $D_{40}.$


\end{enumerate}